\talkslide{beam-rl}{Safety RL}{Timing: provisional}{Safety RL combines several task types, with a judge supplying feedback.}{
\begin{tikzpicture}[remember picture,overlay,
 every node/.style={font=\normalsize,text=Ink,align=center,inner sep=.14cm},
 block/.style={draw=RLBlue,fill=RLBlueTint,rounded corners=5pt,line width=1pt,
  minimum width=3.4cm,minimum height=1.25cm},
 flow/.style={-{Latex[length=2.3mm]},draw=Ink,line width=1pt}]
\definecolor{RLBlue}{HTML}{4777A4}
\definecolor{RLBlueTint}{HTML}{E5EFF8}
\definecolor{RLSage}{HTML}{65814B}
\definecolor{RLSageTint}{HTML}{EDF3E5}
\definecolor{RLWarmGray}{HTML}{A59B8A}
\definecolor{RLWarmTint}{HTML}{F5F2EC}
\begin{scope}[shift={(current page.south west)}]
 \node[block] (policy) at (3.6,6.0) {Policy $\pi_\theta$};
 \node[block,draw=RLWarmGray,fill=RLWarmTint] (environment) at (12.0,6.0)
  {Task / environment};
 \node[block] (update) at (3.6,2.75) {RL update};
 \node[block,draw=RLSage,fill=RLSageTint]
  (judge) at (12.0,2.75) {Safety judge};
 \draw[flow] ([yshift=.3cm]policy.east) -- node[above] {Action}
  ([yshift=.3cm]environment.west);
 \draw[flow] ([yshift=-.3cm]environment.west) -- node[below] {Observation}
  ([yshift=-.3cm]policy.east);
 \draw[flow] (environment.south) -- node[right,align=left]
  {Response\\and context} (judge.north);
 \draw[flow] (judge.west) -- node[above] {Reward} (update.east);
 \draw[flow] (update.north) -- (policy.south);
\end{scope}
\end{tikzpicture}

}
\speak{With RL, the model generates its own response, and a judge supplies feedback for the update. SFT increases the likelihood of desired responses but does not guarantee that every sampled response follows the safety policy. RL trains on the current model's actual responses; it does not guarantee universal compliance either. Safety RL combines direct requests, conversations, jailbreaks, and agentic tool use. The training reward distinguishes harmful from harmless requests: answering a harmless request passes, while refusal or pivoting fails. For a harmful request, the rule reverses. The judge receives previous turns, the original jailbreak seed, or the unsafe agentic action as appropriate. The later evaluation rule additionally includes clarification and incoherent responses. The task/environment produces observations and receives the policy’s actions. This can be a single-turn response or an agentic interaction; it does not imply a learned world model. The safety judge labels behavior, and the task-dependent scoring rule maps that label to a scalar reward for the RL update. Safety data is mixed with the other tasks during the same training run, not added as a separate phase at the end. Beam used a static collection of jailbreaks, not a jointly trained attacker. Source: \href{https://reflection.ai/documents/beam-model-tech-report}{Beam technical report}, Safety Post-Training RL.}
